\documentclass[a4paper,11pt,twoside]{article}
\usepackage[utf8]{inputenc}
\usepackage[T1]{fontenc}
\usepackage[french]{babel}
\usepackage{import}
\import{C:/Users/lasca/Desktop/Cours de mathématiques/(Modèle de cours et d'exercices)}{preambulestandard.tex}
\import{C:/Users/lasca/Desktop/Cours de mathématiques/(Modèle de cours et d'exercices)}{environnementspersonnels.tex}
\import{C:/Users/lasca/Desktop/Cours de mathématiques/(Modèle de cours et d'exercices)}{commandespersonnelles.tex}

\newcommand{\Dev}[2]{
\begin{tcolorbox}[colframe=indigochinois!55!black,colback=indigochinois!5!white,title=Développement. \theme{#1},
shadow={-1.5mm}{-1.5mm}{0mm}{fill=black,
opacity=0.25},
shadow={1.5mm}{-1.5mm}{0mm}{fill=blue,
opacity=0.25}]
#2
\end{tcolorbox}\vspace{0.4cm}}

\usepackage{refcount}
\newcounter{itemcntr}
\makeatletter
\NewDocumentEnvironment{adherencedev}{m}{%
  % #1 is the label for the environment
  \setcounter{itemcntr}{0}%
  % before \item, step the counter if at level 1
  \AddToHook{cmd/item/before}[itemcount]{%
    \ifnum\inteval{\@itemdepth+\@enumdepth}=1 \stepcounter{itemcntr}\fi
  }%
  \begin{itemize}[label=$\star$]%
}{%
  \end{itemize}%
  % remove the hook
  \RemoveFromHook{cmd/item/before}[itemcount]%
  % set the label
  \addtocounter{itemcntr}{-1}%
  \refstepcounter{itemcntr}\label{#1}%
}
\makeatother
\newcommand{\numitems}[1]{\getrefnumber{#1}}

%Pour les nbdp :
\makeatletter
\addtolength{\skip\footins}{-9pt}
\makeatother

\newcommand{\showcorrection}{Commenter cette ligne pour faire disparaître les corrections des exercices}

\sectionfont{\fontsize{13}{15}\selectfont}
\setlength\parindent{0pt}
\linespread{0.9}

\title{\huge Un peu plus de 40 développements gouleyants... \\ pour se tenir à jour}
\date{\today}
\author{Louis Lascaud}

\begin{document}

%Commandes de dernière minute
\makeatletter
\let\ref\@refstar
\makeatother

\thispagestyle{empty}

\maketitle

\vspace{1.4 cm}

\tableofcontents
\addtocontents{toc}{\protect\thispagestyle{empty}}
\vspace{1 cm}

Cette feuille complète le documents « 40 développements gouleyants » en rajoutant des développements et en actualisant le couplage pour qu'il corresponde à l'évolution du programme pour la session en cours. Le nouveau couplage apparaît dans le tableau suivant : \url{https://louis-lascaud.github.io/l-lscd-mathematiques/Docs/depotprive/New%20couplage.xlsx} à télécharger au format Excel. \\ \vspace{0.3 cm}

\begin{large}
\gras{Statistiques} \\
\end{large}

\begin{EnumStarpar}
\item 20 développements d'algèbre, 20 développements d'analyse
\item Par domaines :
\begin{center}
\begin{small}
\begin{tabular}{|c|c|c|c|} \hline
théorie des groupes & 8 & topologie & 5 \\ \hline
algèbre commutative & 3 & analyse réelle & 6 \\ \hline
théorie de Galois & 1 & analyse fonctionnelle & 5 \\ \hline
arithmétique & 4 & calcul différentiel & 4 \\ \hline
algèbre linéaire & 6 & théorie de Fourier & 2 \\ \hline
algèbre bilinéaire & 5 & analyse complexe & 3 \\ \hline
géométrie & 4 & probabilités et stats & 4 \\ \hline
géométrie algébrique & 3 & analyse numérique & 2 \\ \hline
\end{tabular}
\end{small}
\end{center}
\vspace{0.3 cm}

\item 12 développements « mixtes » : Théorème de Lie-Kolchin, Dénombrement sur des permutations aléatoires, L'exponentielle réalise [...], Ellipsoïde de John \& Loewner, Théorème de Kakutani, Théorème de Cartan-von Neumann, Ordres moyens de fonctions arithmétiques, Théorème de Müntz, Théorème de Banach-Alaoglu, Autour de la loi zêta, Théorème de Perron-Frobenius, Méthode de Tchebychev.
\item $r = 3, 4$ de recasage moyen dans le couplage adopté. (Nombre de leçons en \num{2026} : $68 = 34 \times 2$, soit $136$ développements à sélectionner.) Pire recasage : Théorème de Sturm (1 leçon). Meilleurs recasages : Théorème de Burnside (7 leçons), $SO(3) = PSU(2)$, théorème de Lie-Kolchin (6 leçons). Encore une fois vu le poids de certains dévs en recasages on peut tout à fait faire autrement. Typiquement, les développements sont 10, 15, 23, 29, 35 ou 40 sont dispensables.
\item Pour prévoir l'évolution de la liste des leçons, qui peut changer d'une année sur l'autre, voici un lien vers des nouveaux développements, avec le couplage mis à jour en conséquence pour la session courante :
\item Plusieurs développement sont liés entre eux, soit que des notions se recoupent, soit qu'ils utilisent les mêmes objets :
\begin{EnumTriangle}
\item 1, 18 et 38 utilisent de façon cruciale le cas d'égalité de l'inégalité triangulaire.
\end{EnumTriangle}
\end{EnumStarpar}
\vspace{0.8 cm}

\newpage

\section{Théorème « $p^aq^b$ » de Burnside}

Préambule. \\

\gras{Références :} réfs. \\

\gras{Recasages :} leçons. \\

\Dev{Titre du dév}{Énoncé.}

\gras{Dessin :} item \\

\gras{Couleur :} item \\

\gras{Pédagogie :} texte. \\

\gras{Adhérence du développement (\numitems{dev1}) :}
\begin{adherencedev}{dev1}
\item Item
\end{adherencedev}

\section*{Bibliographie pratique mise à jour}
\label{sec:biblio}
\addcontentsline{toc}{section}{\nameref{sec:biblio}}

[Bibliothèque entièrement remodifiée]

\end{document}